% meta.tex
%
% Aetf <aetf@unlimitedcodeworks.xyz>
% Copyright 2016 Aetf <aetf@unlimitedcodeworks.xyz>
%
% multiple1902 <multiple1902@gmail.com>
% Copyright 2011~2012, multiple1902 (Weisi Dai)
%
% Project Home: https://github.com/Aetf/xjtuthesis
%
% It is strongly recommended that you read documentations located at
%   https://github.com/Aetf/xjtuthesis/wiki
% in advance of your compilation if you have not read them before.
%
% This work may be distributed and/or modified under the
% conditions of the LaTeX Project Public License, either version 1.3
% of this license or (at your option) any later version.
% The latest version of this license is in
%   http://www.latex-project.org/lppl.txt
% and version 1.3 or later is part of all distributions of LaTeX
% version 2005/12/01 or later.
%
% This work has the LPPL maintenance status `maintained'.
%
% The Current Maintainer of this work is Aetf.
%

% 标题，中文
\ctitle{基于信息交互的多智能体强化学习算法研究}

% 作者，中文
\cauthor{陶泽睿}

% 学科，中文，本科生不需要
\csubject{}

% 学院，专业，班级，中文，本科毕业论文需要
\ccollege{计算机科学与技术学院}
\cmajor{计算机科学与技术}
\cclass{计算机2204}

% 设计所在单位，中文，本科毕业论文需要
\corgnization{西安交通大学}

% 学号
\stunum{2204112883}

% 导师姓名，中文
\csupervisor{蔺杰}

% 关键词，中文。用全角分号「；」分割
% 研究生的应首先从《汉语主题词表》中摘选
\ckeywords{多智能体强化学习；显式通信；集中式训练分布式执行；MAPPO；动作层介入}

% 提交日期
\cproddate{{\color{black}\the\year} 年\ {\color{black}\twodigits{\the\month}} 月}

% 论文类型，中文，本科生不需要
% 从理论研究、应用基础、应用研究、研究报告、软件开发、设计报告、案例分析、调研报告、其它中选择
\ctype{}

% 论文标题，英文
\etitle{Research on Multi-Agent Reinforcement Learning Algorithms with Information Exchange}

% 作者姓名，英文
\eauthor{Zerui Tao}

% 学科，英文，本科生不需要
\esubject{}

% 导师姓名，英文
\esupervisor{Jie Lin}

% 关键词，英文。用半角分号和一个半角空格「; 」分割
\ekeywords{multi-agent reinforcement learning; explicit communication; centralized training and decentralized execution; MAPPO; action-level intervention}

% 学科门类，英文
% 从Philosophy（哲学）、Economics（经济学）、Law（法学）、Education（教育学）、Arts（文学）、
%   Science（理学）、Engineering Science（工学）、Medicine（医学）、Management Science（管理学）中选择 
\ecate{Engineering Science}

% 提交日期，英文，本科生不需要
% 应当和 cproddate 保持一致
\eproddate{\monthname{\month}\ \the\year}

% 论文类型，英文，本科生不需要
% 从Theoretical Research（理论研究）、Application Fundamentals（应用基础）、Applied Research（应用研究）、
%   Research Report（研究报告）、Software Development（软件开发）、Design Report（设计报告）、
%   Case Study（案例分析）、Investigation Report（调研报告）、其它（Other）中选择
\etype{}

% 摘要，中文。段间空行
\cabstract{
多智能体强化学习在协作博弈、群体控制和分布式决策等任务中具有重要应用价值。在部分可观测环境中，单个智能体仅依赖局部观测往往难以形成稳定协同，因此如何在集中式训练、分布式执行范式下，以较低通信代价实现有效而稳定的显式通信，成为当前研究中的关键问题。

本文以 MAPPO 为基础框架，围绕“通信如何稳定进入动作决策”展开研究。针对冷启动通信易放大训练噪声、消息语义不清晰以及通信难以真实作用于动作层等问题，本文提出两阶段训练范式：先训练稳定的无通信主干策略，再从同一检查点出发进行通信热启动微调。在此基础上，本文构建统一的轻量残差通信框架，将通信限制在任务相关的动作子空间中，并依次设计选择性路由、沉默预算、动作意图共享、队友冲突边际约束和不确定性融合暴露机制，最终形成稳定通信介入方案。同时，本文建立了包含执行期通信代价、通信隔离评估和队友建议相对本地决策影响诊断在内的评价体系。

实验在原型环境 \texttt{join1}、SMAC 主实验地图 \texttt{5m\_vs\_6m} 以及补充地图 \texttt{MMM2}、\texttt{2c\_vs\_64zg} 上进行。结果表明，通信并非天然有效，缺乏结构约束的冷启动通信可能直接破坏训练，而两阶段训练能够显著提高通信实验的稳定性与可解释性；在 \texttt{5m\_vs\_6m} 上，本文方案实现了低通信代价下的稳定动作层介入，其代表性检查点的总执行期通信预算约为 $21.36$，仅为仓库内 MAIC 实现粗粒度预算的 $44.5\%$，并在通信隔离评估中取得了相对禁用攻击通信模式约 $4.7\%$ 的代表性执行期增益；在 \texttt{MMM2} 与 \texttt{2c\_vs\_64zg} 上，本文方案相对于同预算无通信续跑也表现出一定的正向迁移收益。本文研究表明，协作型多智能体显式通信的关键不在于扩大消息容量，而在于构造受约束、可解释且与动作子空间对齐的低成本介入机制。当前工作已经证明通信可以在较低预算下稳定进入动作层并产生可测量收益，但其对本地决策边界的跨越仍未完全解决。
}

% 摘要，英文。段间空行
\eabstract{
Multi-agent reinforcement learning plays a vital role in cooperative games, swarm control, and distributed decision-making, yet partial observability makes stable coordination particularly challenging. A central problem is how to achieve effective yet low-cost explicit communication under centralized training with decentralized execution.

This thesis studies this problem on a MAPPO backbone, focusing on stable communication intervention in action decisions. To address cold-start instability, ambiguous message semantics, and weak action-level influence, a two-phase paradigm is adopted: first train a stable no-communication backbone, then warm-start communication fine-tuning from the same checkpoint. A unified lightweight residual communication framework is then constructed, restricting communication to task-relevant action subspaces. Selective routing, silence-budget regularization, action-intention sharing, teammate-conflict margin constraint, and uncertainty fused exposure are sequentially introduced, forming a stable communication intervention scheme, alongside an evaluation framework covering execution-time communication cost, communication isolation, and teammate-versus-local diagnostic analysis.

Experiments on \texttt{join1}, \texttt{5m\_vs\_6m}, \texttt{MMM2}, and \texttt{2c\_vs\_64zg} show that unconstrained cold-start communication can destabilize training, while the two-phase procedure improves stability and interpretability. On \texttt{5m\_vs\_6m}, the scheme achieves stable action-level intervention at low cost: its total execution-time communication budget is $21.36$, merely $44.5\%$ of the repository MAIC implementation, with a $4.7\%$ gain over the no-attack-communication mode in isolation evaluation. Gains over same-budget no-communication continuation are also observed on \texttt{MMM2} and \texttt{2c\_vs\_64zg}. These findings suggest that effective MARL communication hinges not on expanding message capacity, but on constructing constrained, interpretable mechanisms aligned with action subspaces. This work demonstrates stable communication intervention at the action layer under a low budget, while crossing the local decision boundary remains an open challenge.
}
